% Chapter 2

\chapter{Literature Review} % Main chapter title

\label{Chapter2} % For referencing the chapter elsewhere, use \ref{Chapter2}

%----------------------------------------------------------------------------------------

It is not a new phenomenon that the problem of designing schedules that have no conflicts and that are perfect has received much attention through decades of research. This chapter summarizes the current literature on the UTP, categorizes the different types of algorithms developed by researchers for tackling the problem, and discusses the evolution of software architecture in the context of timetable management systems.

%-----------------------------------
\section{The University Timetabling Problem (UTP)}

Formally speaking, academic scheduling is referred to as the problem of allocating a collection of activities (like lectures, examinations, and meetings) to a collection of resources (time periods, rooms, and teaching staff) based on a variety of constraints \cite{schaerf1999survey}. It should be noted that the difficulty of such a task depends on many factors, such as the scale of the institution, the flexibility of the course structure, and the constraints' strictness.

\subsection{Constraint Satisfaction and Complexity}
In computational complexity theory, the University Timetabling Problem is generally classified as NP-hard \cite{even1975complexity}. This means that when the number of variables (classes, rooms, teachers) increases, the time needed to solve the problem increases exponentially, making any exhaustive search algorithm totally impractical.

Constraints within the UTP are universally divided into two categories:
\begin{itemize}
    \item \textbf{Hard Constraints:} These are non-negotiable conditions that must be satisfied for a timetable to be considered feasible. Examples are:
    \begin{itemize}
        \item A room cannot host more than one class simultaneously.
        \item A professor cannot teach two different classes at the same time.
        \item A student section cannot be scheduled for two different lectures simultaneously.
        \item The seating capacity of an assigned room must meet or exceed the number of students enrolled in the class.
    \end{itemize}
    \item \textbf{Soft Constraints:} These represent preferences that enhance the quality of the timetable but are not strictly necessary for its execution. A penalty score is usually assigned when soft constraints are violated, and the goal of optimization algorithms is to minimize this score. Examples include:
    \begin{itemize}
        \item Minimizing the number of free periods (gaps) in a student's daily schedule.
        \item Ensuring faculty members do not have heavily fragmented schedules across the week.
        \item Assigning classes to rooms preferred by the instructor.
        \item Grouping specific core subjects together sequentially.
    \end{itemize}
\end{itemize}

%-----------------------------------
\section{Algorithmic Approaches to Timetabling}

Over the years, researchers have proposed a vast array of techniques to solve the UTP, ranging from classical mathematical modeling to modern meta-heuristic algorithms. 

\subsection{Exact Methods}
The earliest methods concerned at the mathematical level whether Integer Linear Programming (ILP) and Graph Coloring. In a Graph Coloring model, events are represented as nodes, and edges are drawn between events that conflict (e.g. share the same student group, professor).  The objective is to use the fewest possible colours to colour the nodes, where colours indicate distinct time slots, such that no two adjacent nodes are coloured the same colour. Although theoretically valid, exact approaches have significant difficulties due to the high dimension of real-life timetabling problems, which are often unable to find solutions in acceptable computation times when soft constraints are added.

\subsection{Heuristic and Meta-Heuristic Algorithms}
Because exact methods could not deal with large-size problems, heuristics and meta-heuristics became a prominent focus of the scientific community during the late 1990s and 2000s. The methods do not guarantee the absolute best solution, but they are very good at finding “good enough” (feasible and high-quality) solutions in reasonable time.

\begin{itemize}
    \item \textbf{Genetic Algorithms (GA):} Inspired by natural evolution, GAs represent potential timetables as "chromosomes." Through processes of selection, crossover (combining parts of two schedules), and mutation (randomly altering a scheduled event), populations of timetables evolve over generations to minimize constraint violations \cite{burke1996genetic}.
    \item \textbf{Simulated Annealing (SA) and Tabu Search:} These local search techniques iteratively make small changes to a complete schedule. Tabu search uses a memory structure (a "tabu list") to prevent the algorithm from revisiting recently evaluated schedules, thereby avoiding getting trapped in local optima \cite{glover1997tabu}. Simulated Annealing mimics the metallurgical process of annealing, occasionally accepting "worse" solutions early in the process to explore a broader search space before gradually "cooling" down to focus on strict optimization \cite{kirkpatrick1983optimization}.
\end{itemize}

%-----------------------------------
\section{Evolution of Timetable Management Systems}

Although a lot of work has been done by researchers to understand schedule generation mathematically, it cannot be ignored that software engineering for timetable management has advanced considerably, in tandem with advances made in software engineering as a whole.

\subsection{Legacy Systems: Spreadsheets and Desktop Applications}
Timetables were traditionally hand-drawn on large whiteboards and grids. The rise of personal computing gave birth to spreadsheet programs such as Microsoft Excel, which became the standard in the academic world. Spreadsheets provide an immensely versatile user interface that is easy to read and understand. As observed by many educational technology experts, however, spreadsheets are inherently static in nature. They are unable to incorporate relational data models that would allow referential integrity enforcement and detection of logical inconsistencies without relying on error-prone, custom-built macros (e.g., Visual Basic for Applications) \cite{panko1998what}. The problem with sharing an Excel document is the static nature of its data.

Desktop application solutions emerged in the late 1990s and early 2000s as a solution to the drawbacks associated with spreadsheet software. These systems incorporated relational database management systems such as MySQL or Oracle for storing schedule information, providing superior data integrity compared to flat file storage. Nevertheless, due to the desktop-centric nature of such applications, they also had the same distribution problems. The stakeholders involved, namely students and instructors, could not directly communicate with the software but had to resort to printed documents or HTML outputs produced by administrators.

\subsection{Modern Web-Based Systems}
Cloud computing and contemporary web frameworks have greatly influenced the way institutions manage their logistics through software systems. Modern systems employ full-stack web-based architectures, enabling the storage of data centrally while ensuring universal access using web browsers.

With the use of modern front end frameworks such as React, Angular, or Vue.js, highly interactive Single Page Applications (SPAs) have become possible. In contrast to legacy web applications that necessitated full-page refreshes every time an action is performed by a user, SPAs allow for highly fluid user interactions. This is especially important in the case of timetable interfaces since the user needs to be able to change views quickly, for example, switching from "Room view" to "Professor view" and filtering for particular criteria (e.g., only particular students' sections).

In addition, Backend as a Service (BaaS) providers such as Firebase or Supabase have sped up this process considerably. BaaS platforms remove the burden of setting up servers, creating databases, and securing authentications which allows developers to focus on implementing functionality such as conflict detection and heuristic Excel parsers.

%-----------------------------------
\section{Identification of Gaps in Existing Solutions}

While there is an abundance of off-the-shelf schedulers and extensive academic literature on constraints-based scheduling, it seems that there is still a significant gap in many mid-size to large academic institutions, especially those based in developing countries.

1. \textbf{The Ingestion Barrier:} Most commercial software applications demand manual data entry via their interface or importing the data from rigidly defined CSV files. Any institution that maintains a more elaborate, unique Excel format (including merged cells, cell coloring, and embedded notes) will face major difficulties converting its current workflow to such a rigid system. There is an urgent need for heuristic parsers that can read "human-written" spreadsheets.

2. \textbf{Inadequate Building Occupancy Tools:} Most timetabling applications only produce the desired schedule output (student/faculty). However, no tools are available for campus managers to visualize building occupancy. A "Room Finder" utility or "Master Occupancy Map" is not typically offered as part of the core scheduling engine.

3. \textbf{Disconnect between Generation and Management:} The highly scientific solutions provided in academia for scheduling, which are based on Genetic Algorithms, are characterized by terrible interfaces which are not usable by management teams. On the other hand, well-designed commercial products usually present a fixed schedule that is to be viewed, instead of solved.

The literature therefore suggests the following path: pure algorithmic methods devoid of any consideration of human factors, local spreadsheet-based methods, and cloud-based integrated systems.

\begin{table}[htbp]
\centering
\caption{Comparison of Timetabling Methodologies}
\label{tab:methodology_comparison}
\renewcommand{\arraystretch}{1.5}
\begin{tabularx}{\textwidth}{@{} l X X X @{}}
\toprule
\textbf{Feature} & \textbf{Spreadsheets (Legacy)} & \textbf{Pure Algorithms} & \textbf{Modern Cloud Systems} \\
\midrule
\textbf{Conflict Detection} & Manual \& Error-Prone & Fully Automated & Real-time Automated \\
\textbf{Data Accessibility} & Offline / Localized & Varies (often CLI based) & Cloud-synchronized (Web/Mobile) \\
\textbf{Flexibility} & High (free-form data entry) & Low (strict mathematical constraints) & High (heuristics + constraints) \\
\textbf{Integration} & None & Low & High (APIs, Database links) \\
\textbf{User Interface} & Basic Grid & Academic / Non-existent & Role-based Dashboards \\
\bottomrule
\end{tabularx}
\end{table}

The system designed in this thesis is based on this framework, employing the "Modern Cloud System" framework while incorporating a heuristic parser to overcome the limitations of legacy spreadsheet systems, creating an architectural framework that enables seamless automation of constraint solving systems in the future.